% ============================================================================
% Section 2 - Method Justification
% Packages used: booktabs, tabularx
% ============================================================================
\section{Method Justification}
\label{sec:justification}

Each design decision follows from an observation in Assignment~1 (A1). Table~\ref{tab:justification}
gives the chain observation $\rightarrow$ decision $\rightarrow$ expected effect. The paragraphs
after it add the interpretation and the supporting literature. ``F1''--``F8'' refer to the
eight EDA findings of A1.

\begin{table*}[t]
\centering
\small
\renewcommand{\arraystretch}{1.15}
\begin{tabularx}{\textwidth}{@{}p{0.03\textwidth}XXX@{}}
\toprule
 & \textbf{Evidence from Assignment 1} & \textbf{Design decision} & \textbf{Expected effect}\\
\midrule
1 & Frozen XLM-R vectors group usages by inflected form (highest silhouette in 9/10 words), then era, and least by sense; sense silhouette is negative for all modern-only words (F8). & Fine-tune the encoder on same-sense vs.\ different-sense usage pairs with a contrastive loss (Eq.~\ref{eq:loss}). & Distances reflect sense rather than form, so clusters align better with gold senses (1a) and JSD follows real sense change (1b).\\
2 & The WiC-trained XL-LEXEME beat frozen XLM-R layer 8 on 1a (0.760 vs.\ 0.733) and 1b ($\rho$ 0.492 vs.\ 0.324); \citet{periti2024systematic} report the same ordering on eight benchmarks. & Reproduce WiC-style training ourselves on XLM-R base, with target markers, instead of using XL-LEXEME as a black box. & Frozen and fine-tuned encoders share one backbone, so the ablation measures the effect of our training alone.\\
3 & Only ten labelled words exist, so each fold trains on eight (A1 Section~3.6). & Stage A on the organisers' synthetic English set (2{,}193 multi-sense words) before Stage B on the eight Swedish words; word-piece embeddings frozen. & The encoder learns a general, word-independent notion of ``same sense'' instead of memorising eight words; the other test languages keep their vocabulary representations.\\
4 & Era is a near-perfect shortcut: archaic spellings appear in 71.5\% of historical and 0.8\% of modern usages (F4); in A1, the Baseline~B clusters of \textit{herre} followed era (AMI 0.65) rather than sense (AMI 0.17). & Spelling normalisation before encoding, plus era-aware pair sampling: positives across eras, negatives within an era. & Era stops being a cue for ``same sense'', so fewer clusters split one sense at the era boundary.\\
5 & One cluster per word scores 0.756 and the number of gold senses ranges from 3 to 12 (F6); Baseline~A's selection returned one cluster for 9 of 10 words. & Average-linkage clustering with a per-fold cut chosen where same- and different-sense pairs of the validation word separate best, and a minimum cluster size. & The cut is tied to the trained distance scale (margin $m_c$), so a word gets several clusters only when its senses are actually separated.\\
6 & The target-sense share ranges from 5.7\% to 94.6\% across words (F7); A1's rank threshold fixed the positive share, and its tuned threshold collapsed to predicting nearly everything positive. & Absolute cosine threshold to the definition, tuned on the new macro F1 over both classes. & The predicted positive rate can follow the word, and the all-positive solution (0.236 under the new metric) is no longer optimal.\\
\bottomrule
\end{tabularx}
\caption{Method justification: evidence from Assignment 1, the resulting design decision, and its expected effect.}
\label{tab:justification}
\end{table*}

\paragraph{Why a fine-tuned encoder (rows 1--2).} The EDA showed that a general-purpose
masked language model is not trained to decide whether two usages share a meaning:
the frozen representation is organised by surface form and era first (F8). Our interpretation
is that sense information is present but not dominant in the vector geometry,
so a clustering on raw distances recovers forms and eras. The implication is to change the
geometry directly with a pairwise objective, which is what Word-in-Context training does
\citep{pilehvar2019wic,cassotti2023xllexeme}. A bi-encoder is required rather than a
cross-encoder, because clustering needs one vector per usage that can be compared with every
other usage \citep{reimers2019sbert}. We train the encoder ourselves instead of using the
released XL-LEXEME so that the only difference between the frozen and the fine-tuned system is
the training. This keeps the comparison controlled, as planned in the ablation study.

\paragraph{Why synthetic pre-training and frozen embeddings (row 3).} Fine-tuning on eight words
risks learning those words' senses instead of the skill of telling senses apart. The synthetic
set provides sense-labelled usages for more than two thousand words. It is English, so it
teaches the pairwise skill rather than Swedish senses, and the Swedish stage then adapts to the
target data. The test set covers six languages \citep{tahmasebi2026semeval}, so freezing the
word-piece embeddings protects the vocabulary of languages that appear in neither training stage.

\paragraph{Why era-aware sampling (row 4).} The single historical-to-modern step contains most of
the gold change (F2, F6), and it is also where spelling and style change most (F3, F4). If
training pairs from the same era were mostly positive, the encoder could learn ``same era $\Rightarrow$
same sense'', and Subtask~1b would then reward this shortcut at exactly the pairs that matter. Sampling
positives across eras and negatives within an era makes the shortcut unhelpful for the loss.

\paragraph{Why this clustering (row 5).} Sense-based change measures depend on the number of
clusters: K-means with silhouette always returns at least two clusters
\citep{giulianelli2020analysing,rousseeuw1987silhouettes}, which A1 showed to inflate
$\rho$ through granularity alone. Agglomerative clustering with a minimum cluster size can
return one sense, which is correct for near-monosemous words such as \textit{exportera}
(F6). A contrastive objective with margin $m_c$ gives distances a meaning shared across words
(same sense $\approx 0$, different sense $\gtrsim m_c$). A cut chosen on one word can therefore
transfer to another, which percentile-based cuts on frozen vectors cannot do. Clustering all
periods together avoids aligning separate clusterings across ten periods
\citep{periti2024towards}.

\paragraph{Why JSD and APD (Section~\ref{sec:approach-pipeline}).} The official 1b score ranks
JSD between sense distributions \citep{lin1991divergence}, so JSD must be computed from our
clusters. APD, one of the two strongest simple measures on SemEval-2020 Task~1
\citep{schlechtweg2020semeval,kutuzov2020uio}, needs no clustering. \citet{periti2024systematic}
found that form-based measures such as APD rank change better than sense-based ones. Reporting
both therefore separates errors of the representation (APD also fails) from errors of the
clustering (APD succeeds but JSD fails).

\paragraph{Why an absolute threshold for Subtask~2 (row 6).} The imbalance is per word, not per
period (F7). A within-word rank threshold predicts the same share of positives for every word,
which is wrong by construction for \textit{fröken} (5.7\%) and \textit{exportera} (94.6\%).
An absolute similarity threshold can only adapt if similarities are comparable across words;
the contrastive training is meant to make them so.
